\section{Introduction}\label{c:introduction}

The usual two-cylinder covering of a regular tetrahedron is balanced
at a natural threshold: its total base area is half the tetrahedron's
smallest orthogonal projection area. We ask how the axes can vary across
the two families so as to reduce this cost without leaving a gap.
The answer involves two competing effects. Tilting decreases the
perpendicular area of an intercept patch, but neighboring lines then
need a larger patch to keep covering the tetrahedron. We construct
fields for which the first decrease exceeds the second increase.

A cylinder in $\R^3$ is $B+\R u$, where $u$ is a unit vector and
$B\subset u^\perp$ is measurable with finite planar area $|B|$.
For a convex body $K$, define
\[
 A_{\min}(K)=\min_{|u|=1}|\pi_{u^\perp}K|.
\]
Cylinder covering problems have an antecedent in plank coverings. A plank is the
region between two parallel hyperplanes, and its width is the distance
between them. Bang's plank theorem states that the total width of a
finite plank cover of a convex body is at least the body's minimum width;
Bezdek and Litvak study the corresponding question for
cylinders~\cite[Introduction]{BL2009}.
The half-area question asks whether every finite cylinder cover of $K$
has $\sum_i|B_i|\ge A_{\min}(K)/2$. Bezdek records this question,
and the regular tetrahedron's equality example, as due to
Bang~\cite[Problem~3.1]{Bezdek2009}; the same attribution appears in
Bezdek and Litvak~\cite{BL2009}.

There is a stronger directionwise formulation. For a finite family
$\mathcal C=(B_i+\R u_i)_i$, its relative cost is
\[
 \mathcal R(\mathcal C;K)
   =\sum_i\frac{|B_i|}{|\pi_{u_i^\perp}K|}.
\]
For finite cylinder covers, Bezdek and Litvak proved
$\mathcal R\ge1/3$ for arbitrary
three-dimensional convex bodies and $\mathcal R\ge1$ for ellipsoids
\cite[Theorem~3.1]{BL2009}. Bezdek and Khan's
\emph{1-Codimensional Cylinder Covering Conjecture} proposes
$\mathcal R\ge1/2$ in dimension three
\cite[Definition~3 and Conjecture~4.13]{BK2016}.
Since each denominator is at least $A_{\min}(K)$,
\begin{equation}\label{c:relative-comparison}
 \mathcal R(\mathcal C;K)
 \le\frac{\sum_i|B_i|}{A_{\min}(K)}.
\end{equation}
Thus that conjecture implies the half-area bound. The latter question
is also recorded in Verreault's survey~\cite[Question~4.14]{Verreault2026}.

\begin{theorem}\label{c:main}
Every regular tetrahedron $K\subset\R^3$ of positive edge length admits a finite
cylinder cover with compact parallelogram perpendicular bases such that
\[
 \sum_i|B_i|<\tfrac12 A_{\min}(K),
 \qquad
 \sum_i\frac{|B_i|}{|\pi_{u_i^\perp}K|}<\tfrac12.
\]
It also admits finite covers with bounded open measurable bases
satisfying both inequalities.
\end{theorem}

The directionwise conclusion extends to all tetrahedra. The reason is
the affine invariance of relative cylinder cost, noted by Bezdek and
Litvak~\cite[Section~3]{BL2009}.

\begin{corollary}\label{c:affine-tetrahedra}
Every nondegenerate tetrahedron $T\subset\R^3$ admits a finite cylinder
cover $\mathcal C$ with compact parallelogram perpendicular bases and
$\mathcal R(\mathcal C;T)<1/2$. The same conclusion holds with bounded
open measurable bases.
\end{corollary}
\begin{proof}
Translations preserve all the areas in the relative cost, so consider
an invertible linear map $A$. For a unit vector $u$, put
$v=Au/|Au|$ and define the linear map between perpendicular planes
\[
 Q_u=\pi_{v^\perp}A\big|_{u^\perp}:u^\perp\longrightarrow v^\perp.
\]
It is invertible: if $Q_up=0$, then $Ap$ is parallel to $Au$, so
$p$ is parallel to $u$ and belongs to $u^\perp$, hence is zero.
The image of $B+\R u$ has perpendicular base $Q_uB$, and
$\pi_{v^\perp}(AK)=Q_u(\pi_{u^\perp}K)$. The same positive planar
Jacobian therefore multiplies the numerator and denominator of
$|B|/|\pi_{u^\perp}K|$. Thus the relative cost is unchanged, and the
stated base classes are preserved. Every nondegenerate tetrahedron is
an invertible affine image of a regular one, so Theorem~\ref{c:main}
proves the corollary.
\end{proof}

Together with the lower bound of Bezdek and Litvak, this places the
infimum relative cost of a finite cylinder cover of any nondegenerate
tetrahedron in $[1/3,1/2)$. The absolute half-area inequality in
Theorem~\ref{c:main} concerns regular tetrahedra; the affine argument
preserves each directionwise ratio.

The absolute inequality follows already from the rational field in
Section~\ref{c:rat-rational}; \eqref{c:relative-comparison} gives the
relative inequality for that same cover. The subsequent constructions
explain two further ways to preserve the overlap. They have different
local geometry and different behavior at the apex, even though they
lead to the same strict comparison.

Here is the common construction. A point $p=(p_1,p_2)$ in a planar
label set specifies a segment
\[
 \bigl(s,p_1+sV_1(p),h(p_2+sV_2(p))\bigr),\qquad |s|\le L,
\]
where $h>0$ is a fixed physical scale and $V(p)$ is its slope.
Freezing that slope on an intercept patch of parameter area $A$
gives perpendicular area
\[
 \frac{hA}{\sqrt{1+V_1(p)^2+h^2V_2(p)^2}}.
\]
The scalar product of the unit normals of the intercept plane
$\{x=0\}$ and the plane perpendicular to $(1,V_1(p),hV_2(p))$
gives this projection factor. It suggests an integral cost for a continuously
varying family, but does not by itself produce a finite cover.
The companion article \emph{Finite cylinder approximation of ruled
sets}~\cite[Theorem~1.2]{RuledApproximation2026} provides a finite
square-tile realization for a compact label set $D$ and a field $V$
that is $C^1$ on a neighborhood of $D$ and satisfies
$\operatorname{tr}DV=0$ and $-L^{-2}<\det DV\le0$.
This one construction includes opposite nonzero real eigenvalues
within the strict length bound and square-zero differentials with
unrestricted shear. Its weights use the original velocity coordinates,
and physical base area includes the factor $h$. To obtain a prescribed
physical area error $e$, we use weighted tolerance $e/h$ in the theorem.
We check the relevant conditions only after fixing the field's
regularization or cutoff.

The rational field is derived from a stream function constant on the
two walls of a wedge. A positive denominator makes it smooth at the
apex, while wall tangency makes the tilted straight-line maps preserve
the wedge. The labels assigned to a point by the two complementary
families move in opposite directions to first order. A quartic
enlargement of the label domains leaves a nonnegative square as their
second-order overlap surplus. Its area costs less than the projection
gain. We prove coverage both by comparing two failed inverse labels
and, independently, by locating monotone height thresholds.

The singular power field changes the radial homogeneity. Its derivative
is hyperbolic away from the apex; a cutoff is chosen so that the
differential remains small while two tip cylinders cost less than the
quadratic gain. At the limiting exponent, this field becomes the angular
field with square-zero differential. That endpoint is handled without
a small-shear estimate. It cannot be defined at its apex
by continuity. We remove the apex before applying the companion
theorem and then cover the omitted angular segments by a finite
velocity net.

Section~\ref{c:rat-rational} supplies the first complete construction,
including the strict cost comparison and its finite realization.
Section~\ref{c:rat-threshold} gives the independent threshold proof.
Section~\ref{c:rational-dictionary} relates two coordinate forms of the
rational field and applies the distinct logarithmic-cell and product-cell
approximation theorems. Sections~\ref{c:sing-fields} and
\ref{c:nil-section} develop the singular and square-zero constructions.
The reference area is first computed by projecting entering faces.
Appendix~\ref{c:projection-appendix} gives two further explanations,
through the second moments of a normal frame and the derivative of
translated overlap volume.
